\documentclass[letterpaper,11pt]{article}
\usepackage[utf8]{inputenc}
%\usepackage[T2A]{fontenc}
\usepackage{fontawesome5}
\usepackage{latexsym}
\usepackage[empty]{fullpage}
\usepackage{titlesec}
\usepackage{marvosym}
\usepackage[usenames,dvipsnames]{color}
\usepackage{verbatim}
\usepackage{enumitem}
\usepackage[hidelinks]{hyperref}
\usepackage{fancyhdr}
\usepackage[russian]{babel}
\usepackage{tabularx}
\input{glyphtounicode}



% Custom font
\usepackage[default]{lato}

\pagestyle{fancy}
\fancyhf{} % clear all header and footer fields
\fancyfoot{}
\renewcommand{\headrulewidth}{0pt}
\renewcommand{\footrulewidth}{0pt}

% Adjust margins
\addtolength{\oddsidemargin}{-0.5in}
\addtolength{\evensidemargin}{-0.5in}
\addtolength{\textwidth}{1in}
\addtolength{\topmargin}{-.5in}
\addtolength{\textheight}{1.0in}

\urlstyle{same}

\raggedbottom
\raggedright
\setlength{\tabcolsep}{0in}

% Sections formatting
\titleformat{\section}{
  \vspace{-4pt}\scshape\raggedright\large
}{}{0em}{}[\color{black}\titlerule\vspace{-5pt}]

% Ensure that generate pdf is machine readable/ATS parsable
\pdfgentounicode=1




%-------------------------%
% Custom commands
\begin{document}
% Custom commands
\newcommand{\resumeItem}[1]{
  \item\small{
    {#1 \vspace{-2pt}}
  }
}

\newcommand{\resumeSubheading}[4]{
  \vspace{-2pt}\item
    \begin{tabular*}{0.97\textwidth}[t]{l@{\extracolsep{\fill}}r}
      \textbf{#1} & #2 \\
      \textit{\small#3} & \textit{\small #4} \\
    \end{tabular*}\vspace{-7pt}
}

\newcommand{\resumeSubSubheading}[2]{
    \item
    \begin{tabular*}{0.97\textwidth}{l@{\extracolsep{\fill}}r}
      \textit{\small#1} & \textit{\small #2} \\
    \end{tabular*}\vspace{-7pt}
}

\newcommand{\resumeProjectHeading}[2]{
    \item
    \begin{tabular*}{0.97\textwidth}{l@{\extracolsep{\fill}}r}
      \small#1 & #2 \\
    \end{tabular*}\vspace{-7pt}
}

\newcommand{\resumeSubItem}[1]{\resumeItem{#1}\vspace{-4pt}}

\renewcommand\labelitemii{$\vcenter{\hbox{\tiny$\bullet$}}$}

\newcommand{\resumeSubHeadingListStart}{\begin{itemize}[leftmargin=0.15in, label={}]}
\newcommand{\resumeSubHeadingListEnd}{\end{itemize}}
\newcommand{\resumeItemListStart}{\begin{itemize}}
\newcommand{\resumeItemListEnd}{\end{itemize}\vspace{-5pt}}

\definecolor{Black}{RGB}{0, 0, 0}
\newcommand{\seticon}[1]{\textcolor{Black}{\csname #1\endcsname}}

%-------------------------------------------%
%%%%%%  RESUME STARTS HERE  %%%%%

%----------HEADING----------%
\begin{center}
    \textbf{\Huge \scshape Иван Улитин} \\ \vspace{1pt}
    \seticon{faPhone} \ \small +7 (960) 343 63 95 \quad
    \href{https://t.me/qqq_iwi}{\seticon{faTelegram} \underline{@qqq\_iwi}} \ \quad
    \href{mailto:qqqiwi.prog@gmail.com}{\seticon{faEnvelope} \underline{qqqiwi.prog@gmail.com
}} \quad
    %\href{https://www.linkedin.com/in/}{\seticon{faLinkedin} \underline{linkedin.com/firstlast}} \quad
    \href{https://github.com/QQQiwi}{\seticon{faGithub} \underline{github.com/QQQiwi}}
\end{center}

%-----------SKILLS-----------%
\section{Технические навыки}
    \begin{itemize}[leftmargin=0.15in, label={}]
	\small{\item{
		\textbf{Языки}{: Python, SQL, Rust, JavaScript, Dart, Cypher, Java,  Ruby} \\
		\textbf{Технологии}{: PyTorch, Transformers, vLLM, SLURM, Docker, Kubernetes, GitLab, ClearML, Linux, LLM, RAG, LangChain, pandas, Seaborn, Scikit-learn, AWS, GCP, Git, Dagster/Airflow, SQL, NoSQL, S3, Clickhouse, FastAPI} \\
		\textbf{Концепции}{: анализ данных, теория вероятностей и матем. статистика, машинное обучение, глубокое обучение, компьютерное зрение, обработка естественного языка, временные ряды, менеджмент}
	}}
    \end{itemize}

%-----------EXPERIENCE-----------%
\section{Опыт работы}
\resumeSubHeadingListStart

    \resumeSubheading
    {\href{https://www.sberbank.com/}{\underline{Сбербанк}}}{Сентябрь 2025 -- текущий момент}
    {Ведущий эксперт по исследованию данных}{Москва}
    \resumeItemListStart
        \resumeItem{Поднимал и сопровождал inference-инфраструктуру для крупных LLM на GPU-кластерах: DeepSeek v4, GLM 5.1, Qwen 2.5, Qwen 3, Qwen 3.5 различных размерностей от 4B до 400B, включая MoE-модели.}
        \resumeItem{Работал с мультинодовым запуском и системой джобов SLURM: инициализация, конфигурация и поднятие вычислительных задач на кластерах до 16 нод H100 и 8 нод A100.}
        \resumeItem{Эксплуатировал vLLM последних версий: собирал Docker-образы, настраивал параметры сервинга и оптимизировал конфигурации для повышения throughput моделей.}
        \resumeItem{Разворачивал и поддерживал инфраструктурные сервисы для 50--60 исследователей, включая GitLab и ClearML; поднимал reverse-proxy и обеспечивал доступность внутренних сервисов.}
        \resumeItem{Занимался менеджментом инфраструктуры и перераспределением GPU-ресурсов между исследовательскими командами с учетом приоритетов задач и загрузки кластера.}
    \resumeItemListEnd

    \resumeSubheading
    {\href{link_to_company}{\underline{Orb Intelligence}}}{Сентябрь 2024 -- текущий момент}
    {Data Scientist, Менеджер по разметке данных}{Калифорния, удаленная работа}
    \resumeItemListStart
        \resumeItem{Создание, поддержка и оптимизация дата-пайплайнов.}
            \begin{enumerate}
                \item Опробывание технологии AlloyDB от Google;
                \item Оптимизация ELT-пайплайна, обрабатывающего данные в PostgreSQL, S3, Clickhouse;
                \item Настройка системы оповещений в Telegram с помощью сенсоров в Dagster;
                \item Создание ассетов для модернизации и расширения функционала ELT-пайплайна в Dagster;
                \item Деплой в Kubernetes.
            \end{enumerate}
        \resumeItem{Менеджмент и управление процессом разметки данных.}
            \begin{enumerate}
                \item Анализ рынка труда и формулировка вакансии;
                \item Создание тестового задания и собеседование кандидатов;
                \item Обучение команды разметчиков;
                \item Менеджмент команды, формулировка задач. Размер команды под управлением - 8-14 человек.
            \end{enumerate}
        \resumeItem{Эксперименты, проверка гипотез о применении LLM для решения задачи поиска информации по ключевым словам.}
            \begin{enumerate}
                \item Сравнительный анализ и тестирование LLM-агентов с открытым исходным кодом;
                \item Работа с OpenAI API, Gemini API, Perplexity API. Промпт-инженеринг, создание LLM-пайплайнов;
                \item Формулировка задач для команды разметки данных для создания Ground Truth набора данных для экспериментов;
                \item Оценка значений бизнес-метрик, улучшаемых с помощью данных экспериментов.
            \end{enumerate}
        \resumeItem{Эксперимент с решением задачи векторного поиска.}
            \begin{enumerate}
                \item Аугментация данных для получения различных векторных представлений;
                \item Изучение и эксперименты с векторными базами данных Qdrant, Milvus;
                \item Получение результатов со следующими \textbf{метриками: recall@1 = 0.997, recall@10 = 0.998, recall@30 = 0.998, recall@100 = 0.998}.
            \end{enumerate}
    \resumeItemListEnd
    
    \resumeSubheading
    {\href{link_to_company}{\underline{Институт проблем точной механики и управления РАН}}}{Июнь 2022 -- текущий момент}
    {Инженер-исследователь}{Саратов, удаленная работа}
    \resumeItemListStart
        \resumeItem{Создания мобильного приложения и telegram-бота для классификации степени варикоза у пациента.}

            \begin{enumerate}
                \item Парсинг данных с помощью веб-краулера и вебдрайвера (Scrapy и Selenium), а также обработка данных, полученных в результате работы парсера;
                \item Обучение ResNet50 для классификации фотографий с ногами и обычных фото (бинарная классификация);
                \item Обучение ResNet50 и различных моделей-трасформеров (ViT, DeiT) в качестве baseline для классификации степени варикоза на фото (мультиклассовая классификация);
                \item Написание frontend части на языке Dart (мобильное приложение на фреймворке flutter), которая взаимодействует с python-сервером на библиотеке flask;
                \item Осуществление нагрузочного тестирования с помощью JMeter, юнит-тестирование;
                \item Участие в написании научной статьи и выступление на конференциях по результатам проекта: \href{https://www.mdpi.com/2227-7390/10/19/3571}{\textbf{ссылка на статью с результатами}}.
            \end{enumerate}

        \resumeItem{Классификация спутниковых данных в виде мультивариативного временного ряда для определение аврорального километрового радиоизлучения.}

            \begin{enumerate}
                \item Предобработка исходных данных со спутника.
                \item Поиск и обучение baseline-моделей для задачи (Random Forest, kNN, SVM, CatBoost); Feature engineering для временных рядов;
                \item Оценка качества предсказаний моделей;
                \item Написание тезисов по результатам экспериментов с моделями: \href{https://www.elibrary.ru/bgrpzs}{\textbf{ссылка на аннотацию тезисов в РИНЦ}}.
            \end{enumerate}
        
        \resumeItem{Подготовка учебного курса по анализу временных рядов с помощью МО.}

            \begin{enumerate}
                \item Подготовка и использование ARMA, ARIMA, SARIMA и др. алгоритмов для предсказания значения временного ряда;
                \item Применение разных алгоритмов для feature engineering, актуальных для временных рядов (преобразование Фурье, шейплеты и др.);
                \item Подготовка и использование нейросетей для задач предсказания значения, классификации, генерации временных рядов (например, LSTM);
                \item Написание документации и подготовка материалов для студентов: \href{https://github.com/marina-cyber/perm__ml_ai/tree/main/timeseries}{\textbf{ссылка на репозиторий}}.
            \end{enumerate}
        
        \resumeItem{Подготовка курса по компьютерному зрению.}

            \begin{enumerate}
                \item Подготовка и использования различных методов предобработки изображений (фильтр Прюитт, фильтр Собеля и т.д.);
                \item Подготовка и использование нейросетей (в том числе трансформеров) для задач классификации, сегментации изображений, VQA;
                \item Написание документации и подготовка материалов для студентов.
            \end{enumerate}
        
    \resumeItemListEnd
    
    \resumeSubheading
    {\href{link_to_company}{\underline{TapBank}}}{Апрель 2024 -- Март 2025}
    {Программист, Data Scientist}{Москва, удаленная работа}
    \resumeItemListStart
        \resumeItem{Создание чат-бота для ответов на вопросы клиентов по продукту.}
            \begin{enumerate}
                \item Предобработка неоднородных данных документации (текстовые инструкции, примеры кода, ссылки, изображения);
                \item Создание синтетической выборки вопросов/ответов;
                \item Работа с графовыми и векторными базами данных (neo4j, pinecone, chroma, faiss);
                \item Построение LLM-пайплайнов с помощью LangChain;
                \item Эксперименты с построением архитектур взаимодействия с LLM: RAG, RAPTOR, GraphRAG. Метрика оценки качества - Cosine Similarity между ответом LLM и ответом аннотатора: \textbf{0.74};
                \item Деплой на AWS.
            \end{enumerate}
        \resumeItem{Эксперимент с созданием LLM-агента для генерации кода эндпойнтов.}
            \begin{enumerate}
                \item Создание тестовой выборки относительно разных паттернов проектирования;
                \item Zero-shot, Few-shot реализация на основе OpenAI API;
                \item Оценка с помощью Human Evaluation.
            \end{enumerate}
    \resumeItemListEnd

\resumeSubHeadingListEnd

%-----------PROJECTS-----------%
\section{Проекты}
\resumeSubHeadingListStart

    \resumeProjectHeading
    {\href{https://github.com/QQQiwi/diploma}{\textbf{Выявление аномалий в многомерных временных рядах}} $|$ \emph{Pandas, Seaborn, SciPy, PyTS, transformers}}    {Диплом, 2025}
    \resumeItemListStart
        \resumeItem{Изучение ИИ-методов для обнаружения аномалий в мультивариативных временных рядах и выбор подходящих подходов.}
        \resumeItem{Предобработка спутниковых данных — нормализация, очистка от выбросов и приведение к единому формату.}
        \resumeItem{Тестирование классических моделей (TimeSeriesForest, kNN, CatBoost), сравнение их эффективности, настройка гиперпараметров.}
        \resumeItem{Сведение задачи к задаче компьютерного зрения (создание спектрограмм с различными параметрами на основе исходных данных), эксперименты с CV-моделями (ResNet34, Xception, ViT).}
        \resumeItem{Аугментация, настройка моделей и улучшение предобработки для повышения точности.}
        \resumeItem{Подготовка текста, иллюстраций, таблиц и презентации для защиты работы.}
    \resumeItemListEnd

    \resumeProjectHeading
    {\href{project_link}{\textbf{Сервис сбора статистики о вакансиях}} $|$ \emph{FastAPI, PostgreSQL, Selenium/BS4, React, AWS, GCP}}    {Фриланс, 2024}
    \resumeItemListStart
        \resumeItem{Разработка Python-скрипта для парсинга вакансий по регионам с сервисов Avito, HeadHunter.}
        \resumeItem{Создание логики процессинга описания вакансий для определения её важных составляющих элементов (с помощью расстояния Левенштейна, предварительной лемматизации и т.п.).}
        \resumeItem{Создание web-приложения на React и NodeJS для визуализации получаемой статистики.}
        \resumeItem{Докеризация сервисов, деплой на AWS, GCP.}
        \resumeItem{Написание документации, тех. поддержка продукта.}
    \resumeItemListEnd

    \resumeProjectHeading
{\href{project_link}{\textbf{Мультиклассовая классификация мультивариативного временного ряда}} $|$ \emph{PyTorch}}    {{Летняя практика}, 2021}
    \resumeItemListStart
        \resumeItem {В рамках R\&D во время летней практики в компании Grid Dynamics изучал возможность обнаружения аномалий во временных рядах с помощью алгоритмов компьютерного зрения.}
        \resumeItem{Осуществлял аугментацию путем создания алгоритмов генерации похожих данных на основе доменных знаний о продукте (значения датчиков-буйков в топливных резервуарах) с помощью линейной интерполяции.}
        \resumeItem{Проводил сравнительный анализ алгоритмов и подходов для определения 5 видов аномалий. \textbf{Значения метрики оценки - Area Under PR Curve: 0.776, 0.488, 0.294, 0.842, 0.357}.}
    \resumeItemListEnd

%-----------EDUCATION-----------%
\resumeSubHeadingListEnd
\section{Образование}
    \resumeSubHeadingListStart
    \resumeSubheading
    {Саратовский государственный университет}{Сентябрь 2019 -- Февраль 2025}
    {Оконченное высшее с отл., специалитет "Мат. методы защиты информации"}{Саратов}
    \resumeSubHeadingListEnd


\section{Увлечения}
\begin{itemize}[leftmargin=0.15in, label={}]
	\small{\item{
	Фотография/видеомонтаж, блог, пауэрлифтинг, игра на музыкальных инструментах.
}}
\end{itemize}
%-------------------------------------------%
\end{document}